\documentclass[12pt,a4paper]{article}
\usepackage[UTF8]{ctex}
\usepackage{amsmath,amssymb,amsthm}
\usepackage{geometry}
\usepackage{graphicx}
\usepackage{hyperref}
\usepackage{bm}

\geometry{margin=2.5cm}

\title{向量三重积恒等式：$\bm{a} \cdot (\bm{b} \times \bm{c}) = \det([\bm{a}, \bm{b}, \bm{c}])$的证明}
\author{Modern Robotics}
\date{}

\begin{document}

\maketitle

\section{引言}

在向量代数中，有一个重要的恒等式：
\begin{equation}
\bm{a} \cdot (\bm{b} \times \bm{c}) = \det([\bm{a}, \bm{b}, \bm{c}])
\label{eq:main_identity}
\end{equation}

其中$[\bm{a}, \bm{b}, \bm{c}]$表示以$\bm{a}$、$\bm{b}$、$\bm{c}$为列向量的$3 \times 3$矩阵。

这个恒等式在现代机器人学中非常重要，特别是在：
\begin{itemize}
\item 计算体积（平行六面体的体积）
\item 判断三个向量是否共面
\item 推导反对称矩阵的性质
\item 理解叉乘和行列式的关系
\end{itemize}

本文档将详细证明这个恒等式。

\section{符号定义}

设三个3维向量：
\begin{align}
\bm{a} &= \begin{bmatrix} a_1 \\ a_2 \\ a_3 \end{bmatrix} \\
\bm{b} &= \begin{bmatrix} b_1 \\ b_2 \\ b_3 \end{bmatrix} \\
\bm{c} &= \begin{bmatrix} c_1 \\ c_2 \\ c_3 \end{bmatrix}
\end{align}

\section{证明方法1：直接计算}

\subsection{步骤1：计算叉乘}

根据叉乘的定义，计算$\bm{b} \times \bm{c}$：
\begin{align}
\bm{b} \times \bm{c} &= \begin{bmatrix} b_1 \\ b_2 \\ b_3 \end{bmatrix} \times \begin{bmatrix} c_1 \\ c_2 \\ c_3 \end{bmatrix} \\
&= \begin{bmatrix} b_2 c_3 - b_3 c_2 \\ b_3 c_1 - b_1 c_3 \\ b_1 c_2 - b_2 c_1 \end{bmatrix}
\end{align}

\subsection{步骤2：计算点积}

计算$\bm{a} \cdot (\bm{b} \times \bm{c})$：

\begin{align}
\bm{a} \cdot (\bm{b} \times \bm{c}) &= \begin{bmatrix} a_1 \\ a_2 \\ a_3 \end{bmatrix} \cdot \begin{bmatrix} b_2 c_3 - b_3 c_2 \\ b_3 c_1 - b_1 c_3 \\ b_1 c_2 - b_2 c_1 \end{bmatrix} \\
&= a_1(b_2 c_3 - b_3 c_2) + a_2(b_3 c_1 - b_1 c_3) + a_3(b_1 c_2 - b_2 c_1) \\
&= a_1 b_2 c_3 - a_1 b_3 c_2 + a_2 b_3 c_1 - a_2 b_1 c_3 + a_3 b_1 c_2 - a_3 b_2 c_1 \\
&= a_1 b_2 c_3 + a_2 b_3 c_1 + a_3 b_1 c_2 - a_1 b_3 c_2 - a_2 b_1 c_3 - a_3 b_2 c_1
\label{eq:dot_cross_expanded}
\end{align}

\subsection{步骤3：计算行列式}

计算$\det([\bm{a}, \bm{b}, \bm{c}])$。矩阵$[\bm{a}, \bm{b}, \bm{c}]$为：
\begin{equation}
[\bm{a}, \bm{b}, \bm{c}] = \begin{bmatrix} a_1 & b_1 & c_1 \\ a_2 & b_2 & c_2 \\ a_3 & b_3 & c_3 \end{bmatrix}
\end{equation}

使用3×3矩阵的行列式公式（Sarrus法则）：
\begin{align}
\det([\bm{a}, \bm{b}, \bm{c}]) &= \det\begin{bmatrix} a_1 & b_1 & c_1 \\ a_2 & b_2 & c_2 \\ a_3 & b_3 & c_3 \end{bmatrix} \\
&= a_1 b_2 c_3 + b_1 c_2 a_3 + c_1 a_2 b_3 - a_3 b_2 c_1 - b_3 c_2 a_1 - c_3 a_2 b_1 \\
&= a_1 b_2 c_3 + a_3 b_1 c_2 + a_2 b_3 c_1 - a_3 b_2 c_1 - a_1 b_3 c_2 - a_2 b_1 c_3 \\
&= a_1 b_2 c_3 + a_2 b_3 c_1 + a_3 b_1 c_2 - a_1 b_3 c_2 - a_2 b_1 c_3 - a_3 b_2 c_1
\label{eq:determinant_expanded}
\end{align}

\subsection{步骤4：比较结果}

比较式(\ref{eq:dot_cross_expanded})和式(\ref{eq:determinant_expanded})，我们发现它们完全相同！

因此：
\begin{equation}
\bm{a} \cdot (\bm{b} \times \bm{c}) = \det([\bm{a}, \bm{b}, \bm{c}])
\end{equation}

\textbf{证明完成！}

\section{证明方法2：使用行列式的性质}

\subsection{行列式的多重线性性质}

行列式具有多重线性性质：
\begin{enumerate}
\item 对某一列的线性性：$\det([\bm{a}_1, \ldots, \alpha \bm{a}_i + \beta \bm{b}_i, \ldots, \bm{a}_n]) = \alpha \det([\bm{a}_1, \ldots, \bm{a}_i, \ldots, \bm{a}_n]) + \beta \det([\bm{a}_1, \ldots, \bm{b}_i, \ldots, \bm{a}_n])$
\item 反对称性：交换两列，行列式变号
\end{enumerate}

\subsection{使用叉乘的定义}

叉乘$\bm{b} \times \bm{c}$可以写成行列式的形式：
\begin{equation}
\bm{b} \times \bm{c} = \det\begin{bmatrix} \bm{i} & b_1 & c_1 \\ \bm{j} & b_2 & c_2 \\ \bm{k} & b_3 & c_3 \end{bmatrix}
\end{equation}

其中$\bm{i}$、$\bm{j}$、$\bm{k}$是标准基向量。

展开这个"行列式"（实际上是形式上的）：
\begin{equation}
\bm{b} \times \bm{c} = \bm{i}(b_2 c_3 - b_3 c_2) - \bm{j}(b_1 c_3 - b_3 c_1) + \bm{k}(b_1 c_2 - b_2 c_1)
\end{equation}

这与标准叉乘公式一致。

\subsection{点积与行列式的关系}

现在计算$\bm{a} \cdot (\bm{b} \times \bm{c})$：
\begin{align}
\bm{a} \cdot (\bm{b} \times \bm{c}) &= a_1(b_2 c_3 - b_3 c_2) + a_2(b_3 c_1 - b_1 c_3) + a_3(b_1 c_2 - b_2 c_1) \\
&= \det\begin{bmatrix} a_1 & b_1 & c_1 \\ a_2 & b_2 & c_2 \\ a_3 & b_3 & c_3 \end{bmatrix}
\end{align}

这正好是$\det([\bm{a}, \bm{b}, \bm{c}])$！

\section{证明方法3：几何解释}

\subsection{几何意义}

\textbf{$\bm{a} \cdot (\bm{b} \times \bm{c})$的几何意义}：
\begin{itemize}
\item $\bm{b} \times \bm{c}$是一个向量，其长度等于以$\bm{b}$和$\bm{c}$为邻边的平行四边形的面积
\item 方向垂直于$\bm{b}$和$\bm{c}$构成的平面
\item $\bm{a} \cdot (\bm{b} \times \bm{c})$等于$\|\bm{b} \times \bm{c}\|$乘以$\bm{a}$在$\bm{b} \times \bm{c}$方向上的投影长度
\item 这正好等于以$\bm{a}$、$\bm{b}$、$\bm{c}$为邻边的平行六面体的体积
\end{itemize}

\textbf{$\det([\bm{a}, \bm{b}, \bm{c}])$的几何意义}：
\begin{itemize}
\item 行列式的绝对值等于以$\bm{a}$、$\bm{b}$、$\bm{c}$为列向量的矩阵所表示的线性变换对单位立方体的体积缩放因子
\item 对于$3 \times 3$矩阵，行列式等于以三个列向量为邻边的平行六面体的有向体积
\end{itemize}

\textbf{结论}：两者都表示同一个平行六面体的体积，因此相等！

\subsection{符号约定}

\begin{itemize}
\item 如果$\bm{a}$、$\bm{b}$、$\bm{c}$构成右手系，则体积为正
\item 如果$\bm{a}$、$\bm{b}$、$\bm{c}$构成左手系，则体积为负
\item 这与行列式的符号约定一致
\end{itemize}

\section{重要推论}

\subsection{推论1：循环置换不变性}

由于行列式交换两列变号，我们有：
\begin{align}
\bm{a} \cdot (\bm{b} \times \bm{c}) &= \bm{b} \cdot (\bm{c} \times \bm{a}) \\
&= \bm{c} \cdot (\bm{a} \times \bm{b}) \\
&= -\bm{a} \cdot (\bm{c} \times \bm{b}) \\
&= -\bm{b} \cdot (\bm{a} \times \bm{c}) \\
&= -\bm{c} \cdot (\bm{b} \times \bm{a})
\end{align}

\textbf{记忆方法}：循环置换（$\bm{a} \to \bm{b} \to \bm{c} \to \bm{a}$）不改变符号，交换两个向量变号。

\subsection{推论2：共面性判断}

三个向量$\bm{a}$、$\bm{b}$、$\bm{c}$共面（线性相关）当且仅当：
\begin{equation}
\bm{a} \cdot (\bm{b} \times \bm{c}) = \det([\bm{a}, \bm{b}, \bm{c}]) = 0
\end{equation}

\textbf{几何解释}：如果三个向量共面，则它们构成的平行六面体体积为零。

\subsection{推论3：向量三重积的展开}

\textbf{向量三重积}：$\bm{a} \times (\bm{b} \times \bm{c})$

使用恒等式：
\begin{equation}
\bm{a} \times (\bm{b} \times \bm{c}) = (\bm{a} \cdot \bm{c})\bm{b} - (\bm{a} \cdot \bm{b})\bm{c}
\end{equation}

这个公式可以通过行列式恒等式推导出来。

\section{在机器人学中的应用}

\subsection{应用1：计算体积}

\textbf{问题}：计算以三个向量为邻边的平行六面体的体积。

\textbf{方法}：
\begin{equation}
V = |\bm{a} \cdot (\bm{b} \times \bm{c})| = |\det([\bm{a}, \bm{b}, \bm{c}])|
\end{equation}

\subsection{应用2：判断线性无关性}

\textbf{问题}：判断三个向量是否线性无关。

\textbf{方法}：如果$\det([\bm{a}, \bm{b}, \bm{c}]) \neq 0$，则三个向量线性无关。

\subsection{应用3：反对称矩阵的性质}

在推导反对称矩阵的性质时，这个恒等式非常有用。

\textbf{例子}：证明$[\bm{a}][\bm{b}] - [\bm{b}][\bm{a}] = [[\bm{a}]\bm{b}]$

可以使用向量三重积恒等式来证明。

\subsection{应用4：计算力矩}

\textbf{力矩公式}：$\bm{m} = \bm{r} \times \bm{f}$

在某些情况下，需要计算$\bm{n} \cdot (\bm{r} \times \bm{f})$，这可以写成：
\begin{equation}
\bm{n} \cdot (\bm{r} \times \bm{f}) = \det([\bm{n}, \bm{r}, \bm{f}])
\end{equation}

\section{具体例子}

\subsection{例子1：数值计算}

\textbf{问题}：计算$\bm{a} \cdot (\bm{b} \times \bm{c})$，其中
\begin{align}
\bm{a} &= \begin{bmatrix} 1 \\ 2 \\ 3 \end{bmatrix} \\
\bm{b} &= \begin{bmatrix} 4 \\ 5 \\ 6 \end{bmatrix} \\
\bm{c} &= \begin{bmatrix} 7 \\ 8 \\ 9 \end{bmatrix}
\end{align}

\textbf{方法1}：使用点积和叉乘
\begin{align}
\bm{b} \times \bm{c} &= \begin{bmatrix} 5 \cdot 9 - 6 \cdot 8 \\ 6 \cdot 7 - 4 \cdot 9 \\ 4 \cdot 8 - 5 \cdot 7 \end{bmatrix} = \begin{bmatrix} -3 \\ 6 \\ -3 \end{bmatrix} \\
\bm{a} \cdot (\bm{b} \times \bm{c}) &= 1 \cdot (-3) + 2 \cdot 6 + 3 \cdot (-3) = -3 + 12 - 9 = 0
\end{align}

\textbf{方法2}：使用行列式
\begin{align}
\det([\bm{a}, \bm{b}, \bm{c}]) &= \det\begin{bmatrix} 1 & 4 & 7 \\ 2 & 5 & 8 \\ 3 & 6 & 9 \end{bmatrix} \\
&= 1(5 \cdot 9 - 8 \cdot 6) - 4(2 \cdot 9 - 8 \cdot 3) + 7(2 \cdot 6 - 5 \cdot 3) \\
&= 1(-3) - 4(-6) + 7(-3) = -3 + 24 - 21 = 0
\end{align}

\textbf{结果}：两种方法都得到$0$，说明这三个向量共面（线性相关）。

\subsection{例子2：计算体积}

\textbf{问题}：计算以以下向量为邻边的平行六面体的体积：
\begin{align}
\bm{a} &= \begin{bmatrix} 1 \\ 0 \\ 0 \end{bmatrix} \\
\bm{b} &= \begin{bmatrix} 0 \\ 1 \\ 0 \end{bmatrix} \\
\bm{c} &= \begin{bmatrix} 0 \\ 0 \\ 1 \end{bmatrix}
\end{align}

\textbf{计算}：
\begin{align}
\det([\bm{a}, \bm{b}, \bm{c}]) &= \det\begin{bmatrix} 1 & 0 & 0 \\ 0 & 1 & 0 \\ 0 & 0 & 1 \end{bmatrix} = 1
\end{align}

\textbf{结果}：体积为$1$，这是单位立方体，符合预期。

\section{总结}

向量三重积恒等式$\bm{a} \cdot (\bm{b} \times \bm{c}) = \det([\bm{a}, \bm{b}, \bm{c}])$是向量代数中的一个基本恒等式，它：

\begin{enumerate}
\item \textbf{连接了点积、叉乘和行列式}：展示了这三个概念之间的深刻联系
\item \textbf{有清晰的几何意义}：都表示平行六面体的有向体积
\item \textbf{有多种证明方法}：可以从代数、几何等多个角度理解
\item \textbf{有重要的应用}：在机器人学中用于计算、判断线性无关性等
\end{enumerate}

\textbf{记忆技巧}：
\begin{itemize}
\item 左边：点积和叉乘的组合
\item 右边：三个向量组成的矩阵的行列式
\item 两者相等，都表示体积
\end{itemize}

\section{参考文献}

\begin{itemize}
\item Modern Robotics: Mechanics, Planning, and Control
\item 线性代数教材
\item 向量分析教材
\end{itemize}

\end{document}

